\documentclass{article}
\title{GraphAllocBench}
\author{Zhiheng Jiang, Yunzhe Wang, Ryan Marr, Ellen Novoseller, Benjamin T. Files, Volkan Ustun}
\date{MoDeM workshop, IJCAI-ECAI}

\begin{document}
\maketitle

\begin{abstract}
A flexible benchmark for preference-conditioned multi-objective policy learning.
GraphAllocBench provides an environment for multi-objective RL algorithms, used to
test their capacity to manage many conflicting objectives in a complex resource
management scenario.
\end{abstract}

\section{Motivation}

Most RL benchmarks hand the agent a single scalar reward, which quietly assumes the
hard part --- deciding how much one objective is worth in terms of another --- has
already been solved by whoever wrote the reward function.

Real allocation problems do not arrive pre-scalarized. A policy is useful only if it
can be conditioned on a stated preference over objectives at evaluation time.

\section{The benchmark}

\begin{itemize}
  \item \resumeItem{Graph-structured allocation}{Agents and resources form a graph,
  and the policy decides how the resources are assigned across it.}

  \item \resumeItem{Conflicting objectives}{The environment is built so that
  objectives genuinely trade off against one another, rather than being jointly
  maximizable.}

  \item \resumeItem{Preference conditioning}{Policies are conditioned on a
  preference vector, so a single trained policy is evaluated across the whole
  trade-off surface instead of one point on it.}

  \item \resumeItem{Flexibility}{Problem size and objective structure are
  configurable, making the benchmark usable as a stress test rather than a single
  fixed task.}
\end{itemize}

\section{Publication}

Zhiheng Jiang, Yunzhe Wang, \textbf{Ryan Marr}, Ellen Novoseller, Benjamin T. Files,
Volkan Ustun. ``GraphAllocBench: A Flexible Benchmark for Preference-Conditioned
Multi-Objective Policy Learning.'' \textit{Multi-Objective Decision Making (MoDeM)
workshop, IJCAI-ECAI.}

\end{document}
